\chapter{Критерии оценки методов синхронизации}

Для сравнения методов синхронизации состояния используется набор метрик,
отражающих точность поддержания согласованности клиентского состояния с
состоянием сервера.

\begin{enumerate}
    \item Средняя позиционная ошибка (avg\_error), м. 

    Показывает среднее расхождение между клиентским состоянием и корректным
    состоянием сервера. Низкие значения соответствуют высокой точности метода.

    \item 95-ый процентиль ошибки (p95\_error), м.

    Характеризует типичные отклонения в условиях сетевой нестабильности.
    Позволяет оценить стабильность метода, исключив влияние редких выбросов.

    \item Максимальная ошибка (max\_error), м.

    Отражает наибольшее отклонение состояния за весь период измерений.
    Используется для анализа устойчивости к потере пакетов, вариативности задержки
    и ошибкам предсказания.
\end{enumerate}

Для каждого метода метрики измерялись во времени и после усреднялись по всей выборке.
На основе этих данных проводится количественный сравнительный анализ.
